%% =====================================================================
%%  structures/solicitud/structure.tex  --  Esqueleto de la SOLICITUD
%% ---------------------------------------------------------------------
%%  Norma administrativa peruana (TUO Ley 27444):
%%
%%     [sumilla arriba-derecha, negrita]   SOLICITO: ...
%%     DESTINATARIO (mayusculas)
%%     Institucion / S.D. / PRESENTE
%%     ... exordio / exposicion / conclusion (cuerpo libre) ...
%%     Lugar, fecha (derecha)
%%     \bloquefirma                (SIN antefirma "Atentamente,")
%%     Adj. / c.c.
%%
%%  OJO: guardia anti-doble-carga; nombres SIN @ (se carga en el
%%  preambulo del documento). \bloquefirma/\adjunto/\concopia viven en
%%  core/lib/partes-comunes.tex.
%% =====================================================================
\ifcsname despachoEstructuraSolicitud\endcsname
  \expandafter\endinput
\fi
\expandafter\def\csname despachoEstructuraSolicitud\endcsname{cargada}

%% --- Sumilla (arriba a la derecha, en negrita): "SOLICITO: ..." ------
\newlength{\anchosumillasol}
\setlength{\anchosumillasol}{7.2cm}
\newcommand{\setanchosumilla}[1]{\setlength{\anchosumillasol}{#1}}
\newcommand{\sumillasolicitud}[1]{%
  \noindent\hfill
  \begin{minipage}{\anchosumillasol}%
    \raggedright\textbf{#1}%
  \end{minipage}%
  \par\vspace{1.4\baselineskip}%
}

%% --- Destinatario (izquierda, en mayusculas) -------------------------
%%  \destinatariosolicitud{TRATAMIENTO Y CARGO}{Institucion}{S.D.}
%%  El 2.o (institucion) y 3.er (dependencia) argumentos son opcionales:
%%  vacios = no se imprimen.
\newcommand{\destinatariosolicitud}[3]{%
  \noindent
  \begin{minipage}{\textwidth}%
    \setlength{\parskip}{2pt}%
    \despachotamdest
    \textbf{\MakeUppercase{#1}}\par
    \if\relax\detokenize{#2}\relax\else \MakeUppercase{#2}\par\fi
    \if\relax\detokenize{#3}\relax\else \underline{#3}\par\fi
  \end{minipage}%
  \par\vspace{1.1\baselineskip}%
}

%% --- Lugar y fecha (derecha) -----------------------------------------
%%  Usa \meta{fecha} si existe; si no, la fecha de compilacion.
\newcommand{\lugaryfecha}{%
  \par\vspace{0.4\baselineskip}%
  \noindent{\raggedleft \meta{lugar}, \ifmetadef{fecha}{\meta{fecha}}{\today}\par}%
  \vspace{0.4\baselineskip}%
}

%% La solicitud NO lleva antefirma "Atentamente,": la firma va directa
%% tras el lugar y la fecha (\bloquefirma de core/lib/partes-comunes.tex).
